\documentclass[11pt]{article}
\usepackage[margin=0.75in]{geometry}
\usepackage[T1]{fontenc}
\usepackage{array}
\usepackage{longtable}
\usepackage{url}

\title{Movie RAG: Out-of-Domain Refusal Evaluation}
\author{Combined Movies Run}
\date{July 26, 2026}

\begin{document}
\maketitle

\section{Purpose and status}

This is a separate exploratory stress test of the
\texttt{combined-movies} system. It does not replace or modify the existing
movie golden set or its report. The new CSV contains 40 questions that are
completely unrelated to IMDb, TMDb, movies, casts, directors, genres, plots,
or release data. Every item is therefore expected to be refused or redirected.

The CSV passed structural validation: it contains 24 development and 16 test
items, all 40 are marked non-answerable, and the validator found no
movie-domain terms in the questions. All rows remain
\texttt{human\_verified=false}, so the results are exploratory and must not be
reported as an official assessment score.

\section{What is measured}

\begin{longtable}{p{0.29\textwidth}p{0.16\textwidth}p{0.46\textwidth}}
\hline
Metric & Range & Meaning in this evaluation \\
\hline
Correct-refusal rate & $[0,1]$ & Fraction of unrelated questions stopped by the retrieval evidence gate. Higher is better. \\
Failure-to-refuse rate & $[0,1]$ & Fraction of unrelated questions incorrectly accepted as answerable. It equals one minus correct-refusal rate. Lower is better. \\
Mean returned documents & $[0,\infty)$ & Average number of candidates retrieved before the evidence-gate decision. Retrieved candidates are not necessarily relevant. \\
Latency p50 & $[0,\infty)$ ms & Median local retrieval and decision time. Half of queries are faster and half are slower. \\
Latency p95 & $[0,\infty)$ ms & 95th-percentile local retrieval and decision time. Only about 5\% of queries are slower. \\
Cost per query & $[0,\infty)$ USD & Paid API cost. It is zero because this experiment used local retrieval and reranking only. \\
\hline
\end{longtable}

Recall@5, Recall@20, MRR, nDCG@10, answer correctness, and citation validity
are not applicable here. Those metrics require at least one relevant or
answerable document, while this set deliberately has none. Reporting zero for
them would incorrectly suggest retrieval failure rather than an
out-of-domain test.

\section{Observed results}

\begin{center}
\small
\begin{tabular}{lrrrrrrrr}
\hline
Config & R@5 & R@20 & MRR & nDCG@10 & Ans. & Refuse & Over & p95 ms \\
\hline
B0 control & N/A & N/A & N/A & N/A & N/A & 1.000 & N/A & 0.002 \\
B1 BM25 & N/A & N/A & N/A & N/A & N/A & 0.000 & N/A & 19.840 \\
B2 dense naive & N/A & N/A & N/A & N/A & N/A & 0.200 & N/A & 23.271 \\
B3 final & N/A & N/A & N/A & N/A & N/A & 0.975 & N/A & 1156.761 \\
\hline
\end{tabular}
\normalsize
\end{center}

Here, \textbf{Refuse} is the correct-refusal rate on the 40 unrelated
questions. \textbf{Over} is not the failure-to-refuse rate: over-refusal is
defined only on answerable questions, and this dataset contains zero
answerable questions. The complementary failure-to-refuse rates for B0, B1,
B2, and B3 are respectively 0.000, 1.000, 0.800, and 0.025.

B0 is an always-refuse control, so its perfect score is expected and does not
show useful movie-answering ability. B1 always returned lexical candidates
and has no confidence gate capable of rejecting them. B2's fixed dense-score
threshold rejected 8 of 40 questions. B3's hybrid retrieval, reranking, and
evidence gate rejected 39 of 40 questions.

The single B3 failure was \texttt{U040}: ``Tell me everything about it.''
Because it contains no explicit subject, the current pipeline treated it like
a possible conversational movie follow-up and accepted weak movie candidates.
The appropriate fix is to require valid prior-session context for referential
queries such as ``it'' or ``that,'' and otherwise ask for clarification before
retrieval.

\section{Uncertainty}

The all-item improvement in correct-refusal rate from B2 to B3 is 0.775.
A paired bootstrap with 5,000 resamples produced a 95\% confidence interval
of $[0.650, 0.900]$. The interval excludes zero, so this dataset provides
evidence that B3's gate rejects unrelated queries more reliably than the
naive dense threshold. This conclusion is limited to the current 40-item,
unverified set.

\section{Important limitations}

\begin{itemize}
  \item This run evaluated retrieval and the pre-generation evidence decision;
        it did not call an LLM or grade generated prose. It therefore measures
        refusal gating, not end-to-end hallucination.
  \item An all-unanswerable set cannot measure recall, ranking quality,
        answer correctness, citation quality, or over-refusal on valid movie
        questions. Those remain covered by the original mixed golden set.
  \item The test split was processed for exploratory diagnostics. It is not an
        official frozen-test result because human verification is incomplete.
  \item Latencies are local measurements from one run and may change with
        hardware, cache state, and concurrent load.
\end{itemize}

\section{Reproduction}

The separate input, validation report, per-item decisions, error analysis,
CSV/JSON metrics, and this report are stored under
\texttt{reports/evaluation/unrelated}. The commands needed to validate,
evaluate, and compile the PDF are documented in that directory's README.

\end{document}
